\documentclass[12pt]{article}
\usepackage[utf8]{inputenc}
\usepackage[T1]{fontenc}
\usepackage[margin=1in]{geometry}
\usepackage{setspace}
\usepackage{xcolor}
\usepackage[colorlinks=true, urlcolor=blue!60!black, linkcolor=blue!60!black]{hyperref}
\setlength{\parindent}{2em}
\setlength{\parskip}{1.25em}
\renewcommand{\baselinestretch}{1.0}
\begin{document}
\begin{center}
{\LARGE \textit{Grutter v. Bollinger}, 539 U.S. 306 (2003)\par}
\vspace{0.8em}
{\large Case Brief --- Four-Page Edition\par}
\end{center}
\vspace{1em}

\section{Memory Jogger}

A 5--4 Court upheld the University of Michigan Law School's holistic, race-conscious admissions policy under strict scrutiny, holding that the educational benefits of student body diversity are a compelling interest---a framework effectively ended twenty years later in \textit{Students for Fair Admissions v. Harvard} (2023) (\href{https://www.courtlistener.com/opinion/130155/grutter-v-bollinger/}{CourtListener}).

\section{Detailed Case Facts}

The Law School received more than 3,500 applications a year for about 350 seats. Its 1992 policy, written to comply with Justice Powell's opinion in \textit{Regents of the University of California v. Bakke}, required holistic review of each file---GPA, LSAT, essays, recommendations, and ``soft variables''---and reaffirmed a commitment to enrolling a ``critical mass'' of African-American, Hispanic, and Native-American students so they would not feel isolated or like spokespersons for their race. The admissions director consulted ``daily reports'' on the class's racial composition but denied seeking any fixed number.

Barbara Grutter, a white Michigan resident with a 3.8 GPA and a 161 LSAT, was waitlisted and rejected in 1996. She sued under the Equal Protection Clause, Title VI, and \S~1981. At trial, the Law School's expert testified that race-blind admissions would have cut underrepresented minority enrollment from 14.5\% to 4\% of the class.

\section{Procedural History}

After a 15-day bench trial, the district court held that diversity was not a compelling interest and that ``critical mass'' was practically a quota, and it enjoined the use of race. The Sixth Circuit, en banc, reversed 5--4, treating Powell's \textit{Bakke} opinion as binding. The Supreme Court heard the case alongside \textit{Gratz v. Bollinger}, which struck down Michigan's point-based undergraduate system, and affirmed on June 23, 2003.

\section{Judicial Votes}

Affirmed, \textbf{5--4}.
\begin{itemize}
\item \textbf{Majority:} Justice O'Connor, joined by Justices Stevens, Souter, Ginsburg, and Breyer (Scalia and Thomas joined in part).
\item \textbf{Concurrence:} Justice Ginsburg, joined by Justice Breyer.
\item \textbf{Concurring in part and dissenting in part:} Justice Scalia (with Thomas); Justice Thomas (with Scalia in part).
\item \textbf{Dissents:} Chief Justice Rehnquist, joined by Scalia, Kennedy, and Thomas; Justice Kennedy.
\end{itemize}

\section{Holding}

Attaining the educational benefits of a diverse student body is a compelling state interest in public higher education. A race-conscious program is narrowly tailored if it (1) uses race flexibly as a ``plus'' within individualized, holistic review rather than a quota or point system; (2) seriously considers workable race-neutral alternatives; (3) does not unduly harm nonminority applicants; and (4) is limited in duration. The Title VI and \S~1981 claims failed with the equal protection claim. The 25-year statement was framed as an expectation, though Thomas and Scalia treated it as a sunset.

\section{Analysis of Opinions}

\textbf{Majority (O'Connor).} Endorsing Powell's \textit{Bakke} view, O'Connor held that all racial classifications face strict scrutiny but that context matters. Drawing on academic-freedom traditions, the Court deferred to the Law School's educational judgment, and amicus briefs from businesses and retired military leaders showed diversity's benefits were ``not theoretical but real.'' Unlike the mechanical bonus in \textit{Gratz}, critical mass was a flexible goal, as fluctuating minority enrollment (13.5\% to 20.1\%) showed. Narrow tailoring did not require race-neutral alternatives that sacrifice selectivity, but programs must have a ``logical end point.''

\textbf{Ginsburg (concurring).} Given persistent bias and school segregation, ``one may hope, but not firmly forecast'' an end within a generation.

\textbf{Scalia and Thomas.} Scalia ridiculed the ``mystical `critical mass' justification'' and predicted endless litigation. Thomas, quoting Frederick Douglass, argued that only ``pressing public necessity'' justifies racial classifications, that Michigan had no compelling interest in an elite law school, and that preferences stigmatize and ``mismatch'' their beneficiaries.

\textbf{Rehnquist and Kennedy (dissenting).} Rehnquist's tables showed each group's share of admittees tracking its share of applicants, proving racial balancing. Kennedy accepted diversity as compelling but found that daily reports made race decisive at the margin: ``Deference is antithetical to strict scrutiny, not consistent with it.''

\section{Examples: Future Applications}

\begin{enumerate}
\item \textbf{Same-Side (Valid): Holistic Medical School Review.} Treating race as one flexible plus among many, after race-neutral options failed, would be upheld under \textit{Grutter}---though invalid after \textit{SFFA}.
\item \textbf{Same-Side (Valid): Top Ten Percent Plus Holistic Review.} Automatic admission for top graduates, with race one contextual sub-factor for remaining seats and periodic reviews, satisfies \textit{Grutter}.
\item \textbf{Opposite-Side (Invalid): Point Bonuses or Separate Tracks.} A 15-point racial bonus or a separate minority committee fails narrow tailoring, like \textit{Gratz} and \textit{Bakke}.
\item \textbf{Opposite-Side (Invalid): Police Promotions.} A racial plus for promotion to foster ``cross-racial understanding'' fails, because \textit{Grutter}'s deference was tied to academic freedom.
\item \textbf{Fence-Sitter: Race-Neutral Proxies.} A preference for low-income ZIP codes adopted expressly to maintain racial diversity sits between permissible socioeconomic targeting and impermissible racial proxies.
\end{enumerate}

\section{Critique}

\textbf{Colorblind and empirical critiques.} Conservatives faulted the Court's deference to administrators as diluting strict scrutiny. Richard Sander's ``mismatch'' thesis, vigorously contested, argued that large preferences lower grades and bar passage. In \textit{SFFA} (2023), the Court held that unmeasurable diversity goals and programs lacking a logical end point fail strict scrutiny.

\textbf{Progressive critiques.} Derrick Bell's \textit{Diversity's Distractions} argued that grounding affirmative action in diversity rather than remedying subordination served elite institutions and diverted attention from poverty and K--12 inequity. Others observed that \textit{Grutter} and \textit{Gratz} rewarded holistic opacity over transparent criteria. The case pits anticlassification against antisubordination.

\section{Key Quotations}

\begin{quote}
``We expect that 25 years from now, the use of racial preferences will no longer be necessary to further the interest approved today.''
\end{quote}
\noindent --- Justice O'Connor, Opinion of the Court, 539 U.S. at 343.

\begin{quote}
``In order to cultivate a set of leaders with legitimacy in the eyes of the citizenry, it is necessary that the path to leadership be visibly open to talented and qualified individuals of every race and ethnicity.''
\end{quote}
\noindent --- Justice O'Connor, Opinion of the Court, 539 U.S. at 332.

\begin{quote}
``Strict scrutiny is not `strict in theory, but fatal in fact.' ... Context matters when reviewing race-based governmental action under the Equal Protection Clause.''
\end{quote}
\noindent --- Justice O'Connor, Opinion of the Court, 539 U.S. at 326--327.

\end{document}
